\documentclass[11pt]{scrartcl}

\usepackage[sexy]{evan}
\usepackage{import}
\usepackage[table]{xcolor}
\usepackage{xfp}
\usepackage{multicol}

\newcommand{\gad}{\textcolor{yellow}{$\bigstar$}}
\newcommand{\bad}{\textcolor{red}{$\bigstar$}}

\definecolor{dcol0}{HTML}{C8E6C9}
\definecolor{dcol1}{HTML}{D4E9B3}
\definecolor{dcol2}{HTML}{E5ED9A}
\definecolor{dcol3}{HTML}{FFF59D}
\definecolor{dcol4}{HTML}{FFE082}
\definecolor{dcol5}{HTML}{FFCC80}
\definecolor{dcol6}{HTML}{FFAB91}
\definecolor{dcol7}{HTML}{F49890}
\definecolor{dcol8}{HTML}{E57373}
\definecolor{dcol9}{HTML}{D32F2F}

\makeatletter
\newcommand{\getcolorname}[1]{dcol#1}
\makeatother

\newcommand{\dif}[1]{%
    \edef\colorindex{\number\fpeval{floor(#1)}}%
    \edef\fulltext{#1}%
    \colorbox{\getcolorname{\colorindex}}{%
        \ifnum\colorindex>8
            \textbf{\textcolor{white}{\,\fulltext\,}}%
        \else
            \textbf{\textcolor{black}{\,\fulltext\,}}%
        \fi
    }%
}
% Variable para dificultad (inicial 0)
\newcommand{\thmdifficulty}{0}

% Comando para asignar dificultad antes del problema
\newcommand{\problemdiff}[1]{\renewcommand{\thmdifficulty}{#1}}

% Estilo del problema que incluye dificultad antes del título
\declaretheoremstyle[
    headfont=\color{blue!40!black}\normalfont\bfseries,
    headformat={%
      \dif{\thmdifficulty}\quad \NAME~\NUMBER\ifx\relax\EMPTY\relax\else\ \NOTE\fi
    },
    postheadspace=1em,
    spaceabove=8pt,
    spacebelow=8pt,
    bodyfont=\normalfont
]{problemstyle}

    \declaretheorem[style=problemstyle,name=Problema,sibling=theorem]{problema}
    \declaretheorem[style=problemstyle,name=Problema,numbered=no]{problema*}

\definecolor{yellow4}{RGB}{255, 206, 0}

\title {Los GOD}
\subtitle {ebb camp }

\author{ Emmanuel Buenrostro }

\begin{document}
\maketitle
La idea de esta lista es recopilar problemas que se me hagan particularmente bonitos.

\section{Problemas}

\epigraph{Meh. I'll make it. I always do}{  }


\problemdiff{2}
\begin{problem} 
    Sea $ABCD$ un cuadrado y sea $Y$ un punto dentro de la diagonal $AC$ distinto del punto medio de $AC$. La recta perpendicular al segmento $BY$ que pasa por $Y$ corta a la recta $AD$ en $X$ y a la recta $CD$ en $Z$. Muestra que $AX=CZ$. 
\end{problem}


\problemdiff{3}
\begin{problem}
[OMM 2013/1] 
Se escriben los números primos en orden $p_1=2, p_2=3, p_3=5,\ldots$ Encuentra todas las parejas de números enteros positivos $a$ y $b$ con $a-b\geq 2$, tales que $p_a-p_b$ divide al numero entero $2(a-b)$.
\end{problem}



\begin{problem}
[OMM 2020/1]
A un conjunto de 5 enteros positivos, se le llama virtual si el máximo común divisor de cualesquiera 3 de sus elementos es mayor que 1, pero el máximo común divisor de cualesquiera 4 de sus elementos es igual a 1. Demuestra que en cualquier conjunto virtual, el producto de sus 5 elementos tiene al menos 2020 divisores positivos distintos.
\end{problem}



\problemdiff{4}
\begin{problem}
[ORO 2024/3 \gad]
En cada celda de un tablero de $9\times 9$ se escribe un entero tal que entre cualesquiera dos celdas en la misma fila o columna
que tienen el mismo $n$ escrito en ellas hay al menos $n$ celdas entre ellas. ¿Cu\'al es la suma minima posible de los n\'umeros 
del tablero?
\end{problem}


\problemdiff{4}
\begin{problem}[USAMO 2023/1]
    En un triangulo acutangulo $ABC$, sea $M$ el punto medio de $BC$. Sea $P$ el pie de altura de $C$ hacia $AM$. Supon que el circuncirculo de el triangulo $ABP$ intersecta la linea $BC$ en $B$ y $Q$.Sea $N$ el punto medio de $AQ$. Muestra que $NB=NC$.
\end{problem}

\problemdiff{4}
\begin{problem}
Tienes una matriz $n \times n$ de números reales. En una operación puedes negar los valores de una fila o una columna. Prueba que puedes llevar a que todas las filas y columnas tengan una suma no negativa en una cantidad finita de movimientos.    
\end{problem}

\problemdiff{4}
\begin{problem}
    [OMM 2022/5] 
    Sea $n>1$ un entero positivo y sean $d_1< d_2< \dots< d_m$ sus $m$ divisores positivos de manera que $d_1=1$ y $d_m=n$. Lalo escribe los siguientes $2m$ números en un pizarrón:
\[d_1,d_2,\dots,d_m,d_1+d_2,d_2+d_3,\dots,d_{m-1}+d_m,N\]
donde $N$ es un entero positivo. Después Lalo borra los números repetidos (por ejemplo, si un número aparece dos veces, él borrará uno de los dos). Después de esto, Lalo nota que los números en el pizarrón son precisamente la lista completa de divisores positivos de $N$. Encuentra todos los posibles valores del entero positivo $n$.
    
\end{problem}

\problemdiff{4}
\begin{problem}
    [Mexico TST 2025/Diciembre/1.2]
    Sea $ABCD$ un cuadrilatero convexo, supon que existe un punto $P$ en el interior del cuadrado tal que: 
    \[ \angle PAD = \angle ADP = \angle PCB = \angle CBP =\angle CPD\]
    Sea $O$ el circuncentro de $CPD$. Demuestra que $OA=OB$.
\end{problem}

\problemdiff{5}
\begin{problem} [USA TSTST 2023/1 \gad] % creo
    Sea $ABC$ un triangulo con gravicentro $G$. Los puntos $R$ y $S$ estan en los rayos $GB$ y $GC$, tal que 
    $$\angle ABS=\angle ACR=180-\angle BGC$$
    Prueba que $\angle RAS+\angle BAC=\angle BGC$
\end{problem}

\problemdiff{5}
\begin{problem}
    Sea $ABC$ un triangulo acutangulo, sea $D$ el pie de altura desde $C$. La bisectriz de $\angle ABC$ intersecta $CD$ en $E$ e intersecta al circuncirculo $\omega$ de $ADE$ en $F$. Si $\angle ADF=45$ muestra que $CF$ es tangente a $\omega$.
\end{problem}

\problemdiff{5}
\begin{problem} [Centro 2023/3]
    
    Sean $a, b$ y $c$ números reales positivos tales que $a b+b c+c a=1$. Demostrar que
\[
\frac{a^{3}}{a^{2}+3 b^{2}+3 a b+2 b c}+\frac{b^{3}}{b^{2}+3 c^{2}+3 b c+2 c a}+\frac{c^{3}}{c^{2}+3 a^{2}+3 c a+2 a b}>\frac{1}{6\left(a^{2}+b^{2}+c^{2}\right)^{2}} .\]
    \label{OMCC23-3}
\end{problem}


\problemdiff{5}
\begin{problem} [OMM 2023/6]
    Determine todas las funciones $f: \mathbb{Z}^+ \rightarrow \mathbb {Z}^+$ que satisfacen las dos condiciones siguientes para cualesquiera $m,n \in \mathbb{Z}^+$:
    \begin{enumerate}
        \item $f(m+n) \mid f(m)+f(n)$
        \item $f(m)f(n) \mid f(mn)$
    \end{enumerate}
\end{problem}

\problemdiff{5}
\begin{problem}
    [IMO 2025/5]
    Alice and Bazza are playing the inekoalaty game, a two‑player game whose rules depend on a positive real number $\lambda$ which is known to both players. On the $n$th turn of the game (starting with $n=1$) the following happens:
If $n$ is odd, Alice chooses a nonnegative real number $x_n$ such that
\[
    x_1 + x_2 + \cdots + x_n \le \lambda n.
  \]
If $n$ is even, Bazza chooses a nonnegative real number $x_n$ such that
\[
    x_1^2 + x_2^2 + \cdots + x_n^2 \le n.
  \]
If a player cannot choose a suitable $x_n$, the game ends and the other player wins. If the game goes on forever, neither player wins. All chosen numbers are known to both players.

Determine all values of $\lambda$ for which Alice has a winning strategy and all those for which Bazza has a winning strategy.
\end{problem}


\problemdiff{5}

\begin{problem}

    Each point of a three-dimensional space is colored
    with one of two colors such that
    whenever an isosceles triangle $ABC$ with $AB = AC$
    has vertices of the same color $c$ it follows that
    the midpoint of $BC$ also is colored with $c$.
    Prove that there exists a perpendicular square prism
    with all vertices of equal color.

    
\end{problem}


\end{document}